We study how much of supervised adaptation small decision models recover from right / wrong feedback on their own answers. The evaluation covers text decision tasks on two released model families, Vega and Laya, using benchmark gold answers to simulate feedback. The primary setting collects one verdict per training row; an eight-pass setting measures what repeated feedback adds. On DecisionBench at 0.8B, one on-policy pass improves Vega by $\gainSmallPkOneDb$ accuracy points, while reusing one logged sampled verdict per row for eight optimisation passes gives $\gainSmallLogIpsDb$. Eight passes with fresh verdicts give $\gainSmallPkDb$, against $\gainSmallSupEsDb$ for a matched labelled adapter with early stopping. The logged arm remains significantly below repeated feedback at 0.8B, although it also changes the collection policy and training objective. Repeated feedback often identifies the correct answer, and the expected REINFORCE reward gradient weights the supervised gradient by the model's current probability of that answer. Label-prior controls recover much of the gain on typed-decisions and the public task menu, but not on DecisionBench and emotion. On Laya, training the encoder with verdicts improves the released model, while matched labelled training performs better on typed-decisions and emotion; those encoder comparisons have one seed on each of these tasks. Temperature scaling accounts for much of the calibration improvement, and benchmark shortcuts and training-pool overlap limit generalisation. The contribution is a controlled empirical comparison of feedback budgets and matched adaptation methods, not a new estimator or evidence of performance with real users or multimodal inputs. Code and per-row records are released.
